% ---- score2 standard preamble (from: python3 score2_tools.py preamble) ----
\documentclass[10pt]{extarticle}
\usepackage[a4paper,margin=1.5cm]{geometry}
\usepackage[T1]{fontenc}
\usepackage{amsmath,amssymb}
\usepackage{graphicx}
\usepackage[export]{adjustbox}
\usepackage{multicol}
\usepackage{enumitem}
\usepackage{booktabs,array,tabularx}
\setlength{\columnsep}{0.7cm}
\setlength{\parindent}{0pt}
\setlength{\parskip}{3pt}
\pagestyle{empty}
\makeatletter
% figure registry (figdefs.tex is generated by the extract step)
\newcommand{\qfigdef}[4]{\@namedef{qf@file@#1}{#2}\@namedef{qf@w@#1}{#3}\@namedef{qf@wide@#1}{#4}}
\newcommand{\qf@check}[1]{\@ifundefined{qf@file@#1}{\PackageError{score2}{Unknown figure id #1}{}}{}}
\newcommand{\qf@img}[4][]{\includegraphics[width=#2\dimexpr\@nameuse{qf@w@#4}pt\relax,max width=#3,max height=0.45\textheight,#1]{\@nameuse{qf@file@#4}}}
% Full-width questions: finalize wraps every question that contains a wide figure in \qfullbegin...\qfullend,
% so the WHOLE question (stem, figures, options) is set across the page in reading order, then the two columns
% resume. The question's head - its number, stem and first figure (\qfullsplit marks the end of the head) -
% is kept on one page, so a number is never left at a page foot with its figure overleaf; the rest may break
% across pages normally, so no big gaps are left. Without a \qfullsplit, the whole question is kept together
% if it is under 45% of a page. Never write these yourself.
\newif\ifqf@full
\newif\ifqf@boxed
\newsavebox{\qf@box}
\newcommand{\qf@open}{\global\qf@fulltrue\par\global\qf@boxedtrue
  \setbox\qf@box\vbox\bgroup\hsize\textwidth\linewidth\textwidth}
\newcommand{\qf@emit}[1]{\par\egroup\global\qf@boxedfalse
  \ifdim\dimexpr\ht\qf@box+\dp\qf@box\relax<#1\textheight\noindent\usebox{\qf@box}\par\else\unvbox\qf@box\fi}
\newcommand{\qfullsplit}{\ifqf@boxed\qf@emit{0.85}\fi}
\newcommand{\qf@close}{\ifqf@boxed\qf@emit{0.45}\fi\global\qf@fullfalse}
\newcommand{\qfullbegin}{\end{multicols}\qf@open}
\newcommand{\qfullend}{\qf@close\begin{multicols}{2}}
\newcommand{\qfullstart}{\qf@open}
\newcommand{\qfulldone}{\qf@close}
% \qfig[scale]{id}: block figure, centred, directly where it is written (size set by the tool).
% Fallback if a wide figure is not inside a full-width question: break out of the columns in place - no floats.
\newcommand{\qfig}[2][1]{\qf@check{#2}%
  \ifqf@full\qf@blk{#1}{#2}\else
  \ifnum\@nameuse{qf@wide@#2}=1 \expandafter\@firstoftwo\else\expandafter\@secondoftwo\fi
  {\end{multicols}\begin{center}\qf@img{#1}{\textwidth}{#2}\end{center}\begin{multicols}{2}}%
  {\qf@blk{#1}{#2}}\fi}
\newcommand{\qf@blk}[2]{\par\smallskip{\centering\qf@img{#1}{\linewidth}{#2}\par}\smallskip}
% \qfigi[scale]{id}: inline figure (inside option lists / tables / minipages); never breaks columns.
\newcommand{\qfigi}[2][1]{\qf@check{#2}\qf@img[valign=c]{#1}{\linewidth}{#2}}
\makeatother
% structure macros - lint reads these to check completeness and figure placement
\newcommand{\papersec}[1]{\par\bigskip{\centering\Large\bfseries PAPER-#1\par}\medskip}
\newcommand{\partsec}[2]{\par\medskip\noindent{\large\bfseries PART-#1: #2}\par\smallskip}
\newcommand{\sect}[1]{\par\smallskip\noindent\textbf{#1}\par}
\newcommand{\qn}[1]{\par\smallskip\noindent\textbf{#1.}~}
\newcommand{\paraq}[2]{\par\smallskip\noindent\textbf{Paragraph for Questions #1 and #2}\par}
\newenvironment{qoptions}{\begin{itemize}[leftmargin=1.8em,itemsep=1pt,topsep=2pt,label={}]}{\end{itemize}}
\newcommand{\opt}[2]{\item[(#1)] #2}
% figure registry (generated by score2_tools.py finalize)
\qfigdef{p1_phy_s1_q04_1}{JEE Advanced (2021 Series) - Test 15 (12 Jun 2021)-img/p1_phy_s1_q04_1.png}{58.7}{0}
\qfigdef{p1_phy_s1_q07_1}{JEE Advanced (2021 Series) - Test 15 (12 Jun 2021)-img/p1_phy_s1_q07_1.png}{132.3}{0}
\qfigdef{p1_phy_s2_pre_1}{JEE Advanced (2021 Series) - Test 15 (12 Jun 2021)-img/p1_phy_s2_pre_1.png}{75.6}{0}
\qfigdef{p1_phy_s2_pre_2}{JEE Advanced (2021 Series) - Test 15 (12 Jun 2021)-img/p1_phy_s2_pre_2.png}{82.3}{0}
\qfigdef{p1_phy_s2_pre_3}{JEE Advanced (2021 Series) - Test 15 (12 Jun 2021)-img/p1_phy_s2_pre_3.png}{62.9}{0}
\qfigdef{p1_phy_s2_pre_4}{JEE Advanced (2021 Series) - Test 15 (12 Jun 2021)-img/p1_phy_s2_pre_4.png}{61.7}{0}
\qfigdef{p1_phy_s2_pre_5}{JEE Advanced (2021 Series) - Test 15 (12 Jun 2021)-img/p1_phy_s2_pre_5.png}{66.2}{0}
\qfigdef{p1_phy_s2_q03_1}{JEE Advanced (2021 Series) - Test 15 (12 Jun 2021)-img/p1_phy_s2_q03_1.png}{90.8}{0}
\qfigdef{p1_phy_s3_q01_1}{JEE Advanced (2021 Series) - Test 15 (12 Jun 2021)-img/p1_phy_s3_q01_1.png}{64.2}{0}
\qfigdef{p1_phy_s3_q02_1}{JEE Advanced (2021 Series) - Test 15 (12 Jun 2021)-img/p1_phy_s3_q02_1.png}{106.1}{0}
\qfigdef{p1_phy_s3_q03_1}{JEE Advanced (2021 Series) - Test 15 (12 Jun 2021)-img/p1_phy_s3_q03_1.png}{83.0}{0}
\qfigdef{p1_phy_s3_q04_1}{JEE Advanced (2021 Series) - Test 15 (12 Jun 2021)-img/p1_phy_s3_q04_1.png}{77.9}{0}
\qfigdef{p1_phy_s3_q05_1}{JEE Advanced (2021 Series) - Test 15 (12 Jun 2021)-img/p1_phy_s3_q05_1.png}{52.9}{0}
\qfigdef{p1_chem_s1_q02_1}{JEE Advanced (2021 Series) - Test 15 (12 Jun 2021)-img/p1_chem_s1_q02_1.png}{127.7}{0}
\qfigdef{p1_chem_s3_q04_1}{JEE Advanced (2021 Series) - Test 15 (12 Jun 2021)-img/p1_chem_s3_q04_1.png}{103.2}{0}
\qfigdef{p1_chem_s3_q05_1}{JEE Advanced (2021 Series) - Test 15 (12 Jun 2021)-img/p1_chem_s3_q05_1.png}{94.8}{0}
\qfigdef{p2_phy_s1_q01_1}{JEE Advanced (2021 Series) - Test 15 (12 Jun 2021)-img/p2_phy_s1_q01_1.png}{61.5}{0}
\qfigdef{p2_phy_s1_q03_1}{JEE Advanced (2021 Series) - Test 15 (12 Jun 2021)-img/p2_phy_s1_q03_1.png}{86.9}{0}
\qfigdef{p2_phy_s1_q04_1}{JEE Advanced (2021 Series) - Test 15 (12 Jun 2021)-img/p2_phy_s1_q04_1.png}{79.6}{0}
\qfigdef{p2_phy_s1_q05_1}{JEE Advanced (2021 Series) - Test 15 (12 Jun 2021)-img/p2_phy_s1_q05_1.png}{75.1}{0}
\qfigdef{p2_phy_s1_q06_1}{JEE Advanced (2021 Series) - Test 15 (12 Jun 2021)-img/p2_phy_s1_q06_1.png}{70.6}{0}
\qfigdef{p2_phy_s2_q02_1}{JEE Advanced (2021 Series) - Test 15 (12 Jun 2021)-img/p2_phy_s2_q02_1.png}{64.9}{0}
\qfigdef{p2_phy_s2_q03_1}{JEE Advanced (2021 Series) - Test 15 (12 Jun 2021)-img/p2_phy_s2_q03_1.png}{191.7}{0}
\qfigdef{p2_phy_s2_q04_1}{JEE Advanced (2021 Series) - Test 15 (12 Jun 2021)-img/p2_phy_s2_q04_1.png}{107.7}{0}
\qfigdef{p2_phy_s2_q05_1}{JEE Advanced (2021 Series) - Test 15 (12 Jun 2021)-img/p2_phy_s2_q05_1.png}{73.5}{0}
\qfigdef{p2_phy_s2_q06_1}{JEE Advanced (2021 Series) - Test 15 (12 Jun 2021)-img/p2_phy_s2_q06_1.png}{92.8}{0}
\qfigdef{p2_phy_s2_q07_1}{JEE Advanced (2021 Series) - Test 15 (12 Jun 2021)-img/p2_phy_s2_q07_1.png}{83.5}{0}
\qfigdef{p2_phy_s3_para1_2_1}{JEE Advanced (2021 Series) - Test 15 (12 Jun 2021)-img/p2_phy_s3_para1_2_1.png}{245.2}{0}
\qfigdef{p2_phy_s3_para3_4_1}{JEE Advanced (2021 Series) - Test 15 (12 Jun 2021)-img/p2_phy_s3_para3_4_1.png}{154.8}{0}
\qfigdef{p2_chem_s1_q03_1}{JEE Advanced (2021 Series) - Test 15 (12 Jun 2021)-img/p2_chem_s1_q03_1.png}{231.7}{0}
\qfigdef{p2_chem_s1_q03_B}{JEE Advanced (2021 Series) - Test 15 (12 Jun 2021)-img/p2_chem_s1_q03_B.png}{183.0}{0}
\qfigdef{p2_chem_s1_q03_C}{JEE Advanced (2021 Series) - Test 15 (12 Jun 2021)-img/p2_chem_s1_q03_C.png}{131.1}{0}
\qfigdef{p2_chem_s1_q03_D}{JEE Advanced (2021 Series) - Test 15 (12 Jun 2021)-img/p2_chem_s1_q03_D.png}{181.3}{0}
\qfigdef{p2_chem_s1_q05_A}{JEE Advanced (2021 Series) - Test 15 (12 Jun 2021)-img/p2_chem_s1_q05_A.png}{125.6}{0}
\qfigdef{p2_chem_s1_q05_B}{JEE Advanced (2021 Series) - Test 15 (12 Jun 2021)-img/p2_chem_s1_q05_B.png}{183.5}{0}
\qfigdef{p2_chem_s1_q05_C}{JEE Advanced (2021 Series) - Test 15 (12 Jun 2021)-img/p2_chem_s1_q05_C.png}{155.1}{0}
\qfigdef{p2_chem_s1_q05_D}{JEE Advanced (2021 Series) - Test 15 (12 Jun 2021)-img/p2_chem_s1_q05_D.png}{156.3}{0}
\qfigdef{p2_chem_s1_q07_1}{JEE Advanced (2021 Series) - Test 15 (12 Jun 2021)-img/p2_chem_s1_q07_1.png}{123.5}{0}
\qfigdef{p2_chem_s2_q01_1}{JEE Advanced (2021 Series) - Test 15 (12 Jun 2021)-img/p2_chem_s2_q01_1.png}{38.9}{0}
\qfigdef{p2_chem_s2_q02_A}{JEE Advanced (2021 Series) - Test 15 (12 Jun 2021)-img/p2_chem_s2_q02_A.png}{245.2}{0}
\qfigdef{p2_chem_s2_q02_B}{JEE Advanced (2021 Series) - Test 15 (12 Jun 2021)-img/p2_chem_s2_q02_B.png}{184.4}{0}
\qfigdef{p2_chem_s2_q02_C}{JEE Advanced (2021 Series) - Test 15 (12 Jun 2021)-img/p2_chem_s2_q02_C.png}{180.1}{0}
\qfigdef{p2_chem_s2_q02_D}{JEE Advanced (2021 Series) - Test 15 (12 Jun 2021)-img/p2_chem_s2_q02_D.png}{97.9}{0}
\qfigdef{p2_chem_s2_q07_A}{JEE Advanced (2021 Series) - Test 15 (12 Jun 2021)-img/p2_chem_s2_q07_A.png}{155.6}{0}
\qfigdef{p2_chem_s2_q07_B}{JEE Advanced (2021 Series) - Test 15 (12 Jun 2021)-img/p2_chem_s2_q07_B.png}{130.6}{0}
\qfigdef{p2_chem_s2_q07_C}{JEE Advanced (2021 Series) - Test 15 (12 Jun 2021)-img/p2_chem_s2_q07_C.png}{159.4}{0}
\qfigdef{p2_chem_s2_q07_D}{JEE Advanced (2021 Series) - Test 15 (12 Jun 2021)-img/p2_chem_s2_q07_D.png}{157.0}{0}
% ---- body skeleton ----
% \begin{document}
% \begin{center}{\Large\bfseries JEE (Main+Advanced) Score-II --- Test NN (DD Mon YYYY)}\end{center}
% \begin{multicols}{2}
%   \papersec{1}   (only for two-paper tests)
%   \partsec{1}{PHYSICS} \sect{Section-I (i): Single Correct} \qn{1} ... \qfig{p1_phy_q01_1} ...
% \end{multicols}
% \end{document}

\begin{document}
\begin{center}{\Large\bfseries JEE Advanced (2021 Series) --- Test 15 (12 Jun 2021)}\end{center}
\begin{multicols}{2}
\papersec{1}
\partsec{1}{PHYSICS}
\sect{SEC--1}
\qn{1} If electron of hydrogen atom is replaced by another particle of same charge but of double mass then
\begin{qoptions}
\opt{A}{Bohr's radius will increase}
\opt{B}{Ionisation energy of atom will be doubled}
\opt{C}{Speed of new particle in given state will be lesser than electrons speed in same orbit}
\opt{D}{Gap between energy levels will now be doubled}
\end{qoptions}
\qn{2} A metal spherical shell of mass $m$ is filled completely with nonviscous water of mass $m/2$ and is rolled down an inclined plane of height $h$. The total kinetic energy of shell plus water and the velocity of the centre of the shell when it reaches the bottom of the plane is $K$ and $v$ respectively. If water in the shell is freezed and then it is rolled down, the corresponding values of kinetic energy and velocity are $K'$ and $v'$ respectively. Then :
\begin{qoptions}
\opt{A}{$K=\dfrac32mgh$}
\opt{B}{$K'=\dfrac32mgh$}
\opt{C}{$v=v'$}
\opt{D}{$v>v'$}
\end{qoptions}
\qn{3} A spherical planet of radius $R$ has spherically symmetrical distribution of mass density, varying as square of the distance from the centre, from zero at centre to maximum value $\rho_0$ at its surface.
\begin{qoptions}
\opt{A}{The value of escape velocity of a mass $m$ at the surface of planet is $\sqrt{\dfrac{4\pi G\rho_0R^2}{5}}$.}
\opt{B}{The value of acceleration due to gravity `g' varies inside the planet as cube of the distance from centre.}
\opt{C}{The value of escape velocity from surface is same as the escape velocity from another planet of same total mass \& radius but having uniform mass density.}
\opt{D}{The energy required to impart escape velocity to particles of masses `m' \& `2m', at the surface of planet, will be in ratio 1:2.}
\end{qoptions}
\qn{4} Two concentric conducting spherical shells of radii $R$ and $2R$ having charges $q_A$ and $q_B$ and potential $2v$ and $3v/2$ respectively. Now shell $B$ is earthed by closing switch $S$. Let charges on spherical shells $A$ \& $B$ become $q_A'$ \& $q_B'$ respectively, then
\qfig{p1_phy_s1_q04_1}
\begin{qoptions}
\opt{A}{$q_A/q_B=\dfrac12$ (magnitude).}
\opt{B}{$q_A'/q_B'=1$ (magnitude).}
\opt{C}{potential of $A$ after earthing becomes $\dfrac{3v}2$.}
\opt{D}{potential difference between $A$ and $B$ ($v_A-v_B$) after earthing becomes $\dfrac v2$.}
\end{qoptions}
\qn{5} The magnetic flux ($\phi$) linked with a coil depends on time $t$ as $\phi=at^n$, where $a$ is a constant. The induced e.m.f. in the coil is $e$ :
\begin{qoptions}
\opt{A}{If $0<n<1$, $e=0$}
\opt{B}{If $0<n<1$, $e\ne0$ and $|e|$ decreases with time}
\opt{C}{If $n=1$, $e$ is constant}
\opt{D}{If $n>1$, $|e|$ increases with time}
\end{qoptions}
\qn{6} Sound waves from different source are interfering and pressure at some point in space is given as --\\
$P=P_0\sin(200\,\pi t+\pi/2)\cos(3\,\pi t+\pi/3)$
\begin{qoptions}
\opt{A}{The frequency of one source is above 100Hz and other source is below 100Hz.}
\opt{B}{Beat frequency is 3 Hz}
\opt{C}{Beat frequency is 200 Hz}
\opt{D}{Beat frequency can be 3 Hz as well as 200 Hz}
\end{qoptions}
\qn{7} $n$ number of identical equilateral prism are kept in contact as shown in figure. If deviation through a single prism is $\delta$. Then - ($n$, $m$ are integers)
\qfig{p1_phy_s1_q07_1}
\begin{qoptions}
\opt{A}{If $n=2m$, deviation through $n$ prisms is zero.}
\opt{B}{If $n=2m+1$, deviation through system of $n$ prisms is $\delta$.}
\opt{C}{If $n=2m$, deviation through system of $n$ prisms is $\delta$.}
\opt{D}{If $n=2m+1$, deviation through system of $n$ prisms is zero.}
\end{qoptions}
\sect{SEC--2}
\paraq{1}{3}
Following is a circuit containing an ideal battery of emf $E$ in which the elements $X$ \& $Y$ are unknown \& may be a capacitor $C$, a resistor $R$, a battery of emf $E$, or may be their combination as mentioned in each case. The connecting wires are perfectly conducting.
\qfig{p1_phy_s2_pre_1}
The switch $S$ is initially open \& it is closed at $t=0$.
\end{multicols}
\begin{center}\small\setlength{\tabcolsep}{4pt}
\begin{tabular}{|p{0.9cm}|p{4.2cm}|p{0.9cm}|p{4.2cm}|p{0.9cm}|p{3.0cm}|}\hline
 & \textbf{Column-1} & & \textbf{Column-2} & & \textbf{Column-3}\\\hline
(I) & Element $X$ is $R$\newline Element $Y$ is $C$ & (i) & Entire energy (non--zero) supplied by batteries is dissipated as heat & (P) & \qfigi{p1_phy_s2_pre_2}\rule{0pt}{52pt}\\\hline
(II) & Both elements $X$ \& $Y$ are $R$ each & (ii) & Total charge drawn from the battery(s) after long time is $CE/2$ & (Q) & \qfigi{p1_phy_s2_pre_3}\rule{0pt}{52pt}\\\hline
(III) & Element $X$ is $R$, element $Y$ is a parallel combination of $R$ and $C$ & (iii) & Total heat produced in circuit is less than the energy supplied by the battery(s) & (R) & \qfigi{p1_phy_s2_pre_4}\rule{0pt}{52pt}\\\hline
(IV) & Element $X$ is $R$ element $Y$ is battery $E$ joined with same polarity & (iv) & $I=0$ just after $S$ is closed as well as after long time & (S) & \qfigi{p1_phy_s2_pre_5}\rule{0pt}{44pt}\\\hline
\end{tabular}
\end{center}
\begin{multicols}{2}
\qn{1} Of the given 4 situations, which one has maximum value of current just after closing the switch ?
\begin{qoptions}
\opt{A}{I, iii, P}
\opt{B}{IV, i, Q}
\opt{C}{II, iii, P}
\opt{D}{III, i, S}
\end{qoptions}
\qn{2} In which case the charge drawn from the battery(s) during $t=0$ to $t=RC$ is the least ?
\begin{qoptions}
\opt{A}{II, i, S}
\opt{B}{I, iii, P}
\opt{C}{III, ii, Q}
\opt{D}{IV, i, Q}
\end{qoptions}
\qn{3} In which case the current through battery(s) has different non zero values just after closing the switch \& after a long time ?
\begin{qoptions}
\opt{A}{II, i, S}
\opt{B}{I, iv, Q}
\opt{C}{IV, i, Q}
\opt{D}{III, iii, R}
\end{qoptions}
\paraq{4}{6}
A cuboid is filled with liquid of density $\rho_2$ upto height $2h$ \& with liquid of density $\rho_1$, upto height $h$ as shown in the figure and assume pressure at top to be zero.
\qfig{p1_phy_s2_q03_1}
\end{multicols}
\begin{center}\small\setlength{\tabcolsep}{4pt}
\begin{tabular}{|p{0.9cm}|p{4.6cm}|p{0.9cm}|p{3.0cm}|p{0.9cm}|p{3.8cm}|}\hline
 & \textbf{Column-1} & & \textbf{Column-2} & & \textbf{Column-3}\\\hline
(I) & Force on face $ABCD$ due to liquid of density $\rho_1$ & (i) & zero & (P) & Effectively acts at centre of that face\\\hline
(II) & Force on face $ABCD$ due to liquid of density $\rho_2$ & (ii) & $\dfrac{\rho_1gh^2\ell}2$ & (Q) & Effectively acts at $\dfrac h3$ above the bottom of that face\\\hline
(III) & Magnitude of increase in force on face $CDEF$ due to filling liquid of density $\rho_1$ & (iii) & $2\rho_1gh^2\ell$ & (R) & Effectively acts at $\dfrac{4h}3$ below the top of that face\\\hline
(IV) & Force on face $CDEF$ if liquid of density $\rho_2$ only is present & (iv) & $2\rho_2gh^2\ell$ & (S) & Cannot produce a net torque about bottom edge of that face\\\hline
\end{tabular}
\end{center}
\begin{multicols}{2}
\qn{4} Which force has the second highest magnitude (given $\rho_1<\rho_2$) ?
\begin{qoptions}
\opt{A}{I, ii, Q}
\opt{B}{IV, i, Q}
\opt{C}{II, i, S}
\opt{D}{III, iii, P}
\end{qoptions}
\qn{5} Which force produces minimum (Non zero) torque about side $EF$ (given $\rho_1=\rho_2$) ?
\begin{qoptions}
\opt{A}{I, ii, Q}
\opt{B}{II, i, S}
\opt{C}{III, iii, P}
\opt{D}{IV, iii, S}
\end{qoptions}
\qn{6} Which force is equal to the total weight of liquid $\rho_2$ [Given $GE=h$]?
\begin{qoptions}
\opt{A}{I, iii, P}
\opt{B}{II, i, S}
\opt{C}{III, iii, P}
\opt{D}{IV, iv, R}
\end{qoptions}
\sect{SEC--3}
\qn{1} A small putty of mass $m=1$ kg and speed $v_0=4$ m/s, strikes a dumbbell placed on a smooth table as shown in figure. Putty is moving in the plane of the dumbbell. The putty sticks to dumbbell ball which is to be treated as particle. How much thermal energy is produced (in J). Consider rod of dumbbell to be massless.
\qfig{p1_phy_s3_q01_1}
\qn{2} In the figure shown a cart moves on a smooth horizontal surface with a constant velocity $v$ acted upon by an external constant force of magnitude $F$. Sand falls on to the cart with negligible velocity at constant rate $\mu$ kg/s and sticks to the cart. Find $N$ if the power supplied by external agent applying the force $F$ is $\mu v^N$.
\qfig{p1_phy_s3_q02_1}
\qn{3} A torch bulb rated at 4.5 W, 1.5 V is connected as shown in the figure. To make the bulb glow at full intensity, what should be the emf (in volt) of the cell?
\qfig{p1_phy_s3_q03_1}
\qn{4} Given $m_A=30$ kg, $m_B=10$ kg, $m_C=20$ kg. Between $A$ and $B$; $\mu_1=0.3$, between $B$ and $C$; $\mu_2=0.2$ and between $C$ and ground $\mu_3=0.1$. The least horizontal force which can start the motion of any part of the system of three blocks resting upon one another as shown in figure is $F$. Write the value of $F/12$ ($g=10$ m/s$^2$)
\qfig{p1_phy_s3_q04_1}
\qn{5} A circular disc suspended from a torsional wire, oscillates with a time period $T$ when twisted by a small angle $\theta$ from equilibrium and released.
\qfig{p1_phy_s3_q05_1}
If another disc of same material but twice the radius and one fourth thickness is made to oscillate (alone) with same wire, its time period is found to be $nT$. Write the value of $n$.
\partsec{2}{CHEMISTRY}
\sect{SEC--1}
\qn{1} Which of the following statement(s) is/are correct ?
\begin{qoptions}
\opt{A}{Glucose \& Fructose are optical isomers.}
\opt{B}{Fructose can be oxidised by tollens reagent.}
\opt{C}{Monosaccharides containing six membered ring can be called pyranose sugar.}
\opt{D}{Sucrose is a disaccharide and it contains two glycosidic linkages.}
\end{qoptions}
\qn{2} Which of the following statement(s) is/are incorrect ?
\begin{qoptions}
\opt{A}{Bakelite is a polyamide}
\opt{B}{All Biodegradable polymers are natural polymers}
\opt{C}{Azulene \qfigi{p1_chem_s1_q02_1}}
\opt{D}{Sulphonation of $\mathrm{C_6H_6}$, $\mathrm{C_6D_6}$ \& $\mathrm{C_6T_6}$ by Electrophilic aromatic substitution mechanism proceed at same rate.}
\end{qoptions}
\qn{3} Select incorrect statement(s) :
\begin{qoptions}
\opt{A}{pka value of $\mathrm{B(OH)_3}$(aq.) increases on adding glycerol.}
\opt{B}{aqueous solution of $\mathrm{Ca(OH)_2}$ can not be used for removal of temporary hardness of water.}
\opt{C}{Crystals of borax on heating does not produce any neutral oxide.}
\opt{D}{Two layers of solid orthoboric acid are connected together by hydrogen bonds.}
\end{qoptions}
\qn{4} For thermal decomposition of $\mathrm{NH_4NO_3(s)}$ :\\
The correct information(s) is/are -\\
$\mathrm{NH_4NO_3(s)}\xrightarrow{\Delta_rH}\mathrm{N_2O(g)+2H_2O(g)}$ \quad (at $T$ K \& $P$ atm)
\begin{qoptions}
\opt{A}{Heat of reaction is positive}
\opt{B}{Entropy change of the reaction must be $\dfrac{\Delta_rH}{T}$}
\opt{C}{Heat of reaction at constant volume is equal to ($\Delta_rH-3RT$)}
\opt{D}{$\Delta_rG$ of this reaction may be negative at high temperature.}
\end{qoptions}
\qn{5} Select the correct statement(s) regarding following galvanic cell taken at temperature $T$K.\\
$\mathrm{Pt\,|\,H_2(g,\,1\ bar)\,|\,HCl\,(10^{-3}\,M)\,|\,AgCl(s)\,|\,Ag(s)}$\\
Given that : $\dfrac{2.303RT}F=0.06$volt\\
\& $E^\circ_{\mathrm{Ag|AgCl|Cl^-}}=-0.18$volt
\begin{qoptions}
\opt{A}{Cell potential of this cell is 0.54 volt.}
\opt{B}{During discharging, $\mathrm{H_2(g)}$ converts into $\mathrm{H^+}$ (aq.)}
\opt{C}{$\mathrm{H_2(g)}$ is obtained at cathode during working of this galvanic cell.}
\opt{D}{$E^\circ_{\mathrm{Ag^+|Ag}}$ is equal to 0.78 volt if $K_{sp}$ of AgCl is $1\times10^{-10}$.}
\end{qoptions}
\qn{6} Select the correct option(s) about valence Bond Theory regarding co-ordination compounds:
\begin{qoptions}
\opt{A}{It does not explain visible spectra of complexes.}
\opt{B}{It is not able to explain diamagnetic behaviour of $\mathrm{[Co(H_2O)_6]Cl_3}$ even $\mathrm{Co^{3+}}$ is paramagnetic.}
\opt{C}{It is unable to differentiate between weak and strong field ligands.}
\opt{D}{It makes exact predictions regarding the tetrahedral and square planar structures of 4-coordinate complexes.}
\end{qoptions}
\qn{7} Which of the following step(s) is/are involved in extraction of iron from haematite ore ?
\begin{qoptions}
\opt{A}{Gravity separation}
\opt{B}{Reduction by CO}
\opt{C}{Self-reduction}
\opt{D}{Slag formation}
\end{qoptions}
\sect{SEC--2}
\paraq{1}{3}
For given solution in the questions molarity can be treated as molality and solute(s) is/are nonvolatile in nature.\\
Given that : $K_a$ of $\mathrm{CH_3COOH}=10^{-5}$ \quad ; \quad $K_f$ = Molal depression constant\\
\hspace*{3.4em}$K_a$ of $\mathrm{HCOOH}=10^{-4}$ \quad ; \quad $K_b$ = Molal elevation constant
\end{multicols}
\begin{center}\small\setlength{\tabcolsep}{4pt}
\begin{tabular}{|p{0.9cm}|p{3.4cm}|p{0.9cm}|p{5.2cm}|p{0.9cm}|p{3.8cm}|}\hline
\multicolumn{2}{|p{\dimexpr 4.3cm+2\tabcolsep\relax}|}{\centering \textbf{Column - I}\newline\textbf{(Colligative properties)}} & \multicolumn{2}{p{\dimexpr 6.1cm+2\tabcolsep\relax}|}{\centering \textbf{Column - II}\newline\textbf{(One of the product formed during electrolysis at Pt - electrode)}} & \multicolumn{2}{p{\dimexpr 4.7cm+2\tabcolsep\relax}|}{\centering \textbf{Column - III}\newline\textbf{(Nature of aqueous solution at 25$^\circ$C)}}\\\hline
(I) & $\Delta T_b\simeq0.3\,K_b$ & (A) & Metal becomes deposited at cathode & (P) & pH = 8.5\\\hline
(II) & $\pi\simeq0.4\,RT$ & (B) & $\mathrm{H_2}$ gas at cathode & (Q) & pH = 4.7\\\hline
(III) & $\Delta T_f\simeq0.2\,K_f$ & (C) & $\mathrm{H_2}$ gas at anode & (R) & pH $\approx$ 7\\\hline
(IV) & $\pi\simeq0.3\,RT$ & (D) & $\mathrm{O_2}$ gas at anode & (S) & Weakly basic solution\\\hline
\end{tabular}
\end{center}
\begin{multicols}{2}
\qn{1} Which of the following option is correctly matched for 0.1 M HCOOK(aq.)
\begin{qoptions}
\opt{A}{(I),(B),(S)}
\opt{B}{(IV), (B), (S)}
\opt{C}{(III),(D), (Q)}
\opt{D}{(III),(C),(P)}
\end{qoptions}
\qn{2} Select incorrect match for 0.1 m $\mathrm{Ca(NO_3)_2}$ aqueous solution.
\begin{qoptions}
\opt{A}{(I),(B),(R)}
\opt{B}{(III),(A),(P)}
\opt{C}{(IV),(D),(R)}
\opt{D}{(I),(D),(R)}
\end{qoptions}
\qn{3} Select the correct match for aq. solution having 0.1 M $\mathrm{CH_3COONa}$ and 0.2 M $\mathrm{CH_3COOH}$.
\begin{qoptions}
\opt{A}{(II),(B),(Q)}
\opt{B}{(II),(A),(Q)}
\opt{C}{(IV),(A),(S)}
\opt{D}{(I),(C),(R)}
\end{qoptions}
\paraq{4}{6}
\end{multicols}
\begin{center}\small\setlength{\tabcolsep}{4pt}
\begin{tabular}{|p{0.9cm}|p{4.2cm}|p{0.9cm}|p{3.6cm}|p{0.9cm}|p{4.6cm}|}\hline
\multicolumn{2}{|p{\dimexpr 5.1cm+2\tabcolsep\relax}|}{\centering \textbf{Column - I}\newline\textbf{(Name reaction or test)}} & \multicolumn{2}{p{\dimexpr 4.5cm+2\tabcolsep\relax}|}{\centering \textbf{Column - II}\newline\textbf{(Related Reagent or intermediate)}} & \multicolumn{2}{p{\dimexpr 5.5cm+2\tabcolsep\relax}|}{\centering \textbf{Column - III}\newline\textbf{(Miscellaneous)}}\\\hline
(I) & Hoffman bromamide reaction & (i) & Carbene & (P) & Reaction involves upgradation of chain\\\hline
(II) & Reimer tieman reaction & (ii) & $\mathrm{CHCl_3}$ & (Q) & Reaction involves chain degration\\\hline
(III) & Aldol condensation & (iii) & KOH & (R) & Synthesis of amine takes place\\\hline
(IV) & Isocyanide test & (iv) & $\mathrm{Br_2}$ & (S) & Used to test primary amines\\\hline
\end{tabular}
\end{center}
\begin{multicols}{2}
\qn{4} Which of the following option is correctly matched ?
\begin{qoptions}
\opt{A}{(I), (iv), (P)}
\opt{B}{(II), (iv), (Q)}
\opt{C}{(III), (i), (R)}
\opt{D}{(IV), (ii), (S)}
\end{qoptions}
\qn{5} Select incorrect match :
\begin{qoptions}
\opt{A}{(I), (iv), (Q)}
\opt{B}{(II), (i), (R)}
\opt{C}{(III), (iii), (P)}
\opt{D}{(IV), (iii), (S)}
\end{qoptions}
\qn{6} How many reaction/test from column-I can use KOH as reagent ?
\begin{qoptions}
\opt{A}{1}
\opt{B}{2}
\opt{C}{3}
\opt{D}{4}
\end{qoptions}
\sect{SEC--3}
\qn{1} Lime stone($X$) $\xrightarrow{\Delta}$ Solid ($A$) + gas ($B$)\\
Other Binary Compound of Calcium($Y$) $\xrightarrow{\mathrm{H_2O}}$ Solution ($C$) + Gas ($D$)\\
Gas ($B$) turns solution ($C$) milky but does not react with $\mathrm{KMnO_4/H^+}$.\\
Gas ($D$) produce white ppt. in tollen reagent and it also produce benzene when passed through red hot iron tube.\\
The difference of oxidation state of same central atom in ($B$) \& ($D$) is :
\qn{2} How many comparison(s) are incorrect among given ?\\
(i) H $>$ Li (ionisation energy) \quad (ii) Li $>$ Be (size)\\
(iii) Na $>$ Rb (size) \quad (iv) O $>$ N (size)\\
(v) S $>$ P (ionisation energy) \quad (vi) C $>$ O (size)\\
(vii) B $>$ Al (size) \quad (viii) F $>$ Cl (ionisation energy)\\
(ix) N $>$ O (ionisation energy)
\qn{3} Number of non-radioactive alkali metals forming superoxide as major product on heating with excess $\mathrm{O_2}=x$,\\
Number of II-A metals of periodic table whose hydrated halides suffer hydrolysis on heating $=y$.\\
Number of non radio active alkali metals which dissolve in liquid ammonia and produce blue colour solution $=u$\\
The value of $(x+y-u)$ would be\ldots.
\qn{4} Total number of stereo isomers corresponding to structure
\qfig{p1_chem_s3_q04_1}
\qn{5} The formation of dextrorotatory isomer $A$ from its levorotatory isomer $B$ in water is reversible reaction having 1st order kinetics in both direction. Racemic mixture is the equilibrium product if reaction starts from either $A$ or $B$. Analysis of polarimeter experiment showed that 10\% of $B$ was converted into $A$ in 230.3 sec at 25$^\circ$C.\\
($\log5=0.7$, $\log3=0.48$ and $\log2=0.3$)\\
Based on given information select number of correct statement(s) among following.\\
(A) Rate constant of forward reaction at 25$^\circ$C is $5\times10^{-4}\ \mathrm{s^{-1}}$\\
(B) Rate constant of backward reaction at 25$^\circ$C is $5\times10^{-4}\ \mathrm{s^{-1}}$\\
(C) Activation energy of forward reaction is equal to activation energy of backward reaction.\\
(D) Equilibrium constant ($K$) at 25$^\circ$C of $B\rightleftharpoons A$ is 1.\\
(E) Equilibrium constant increases by increasing temperature.\\
(F) Molar entropy of $A$ and $B$ at 25$^\circ$C must be same.\\
(G) $\Delta_rH$ of the process is zero.\\
(H) $\Delta_rG^\circ$ of the process is zero.\\
(I) $\Delta_rG$ of the process \qfigi{p1_chem_s3_q05_1}
\partsec{3}{MATHEMATICS}
\sect{SEC--1}
\qn{1} Let $\vec a,\vec b,\vec c$ be vectors such that $\vec b\cdot\vec c=2\vec a\cdot\vec a=3\vec b\cdot\vec b=\vec a\cdot\vec b\times\vec c=1$ and $\vec a\cdot\vec b=\vec a\cdot\vec c=0$. Now if $\vec a\times\vec b$ is expressed as $x\vec a+y\vec b+z\vec c$, then
\begin{qoptions}
\opt{A}{$x<0$}
\opt{B}{$y<0$}
\opt{C}{$z<0$}
\opt{D}{$x+y+z<0$}
\end{qoptions}
\qn{2} Let $f$ be a twice-differentiable real valued function satisfying the equation :\\
$x+f(x)\,f\,'(x)+f\,'(x)f\,''(x)=0$, then
\begin{qoptions}
\opt{A}{$\lim\limits_{x\to0}xf(x)=-1$}
\opt{B}{$\lim\limits_{x\to0}\dfrac{f\,'(x)}{\ln|x|}=0$}
\opt{C}{$\lim\limits_{x\to0}xf(x)f\,'(x)=1$}
\opt{D}{$\lim\limits_{x\to0}\dfrac{f(x)f\,'(x)}{\ln|x|}=0$}
\end{qoptions}
\qn{3} Let the conics $x^2=xy+1$ and $xy=x^2+1$ have eccentricities $e_1$ and $e_2$ respectively then
\begin{qoptions}
\opt{A}{$e_1>e_2$}
\opt{B}{$e_1<e_2$}
\opt{C}{$e_1\cdot e_2=2$}
\opt{D}{$e_1^{-2}+e_2^{-2}=1$}
\end{qoptions}
\qn{4} Let $A=[a_{ij}]_{3\times3}$ be a matrix where $a_{ij}\in\{0,1\}\ \forall\ 1\le i,j\le3$. Let $s_i=\displaystyle\sum_{k=1}^3a_{ik}$ and $t_j=\displaystyle\sum_{k=1}^3a_{kj}$ then the probability that the product ($s_1s_2s_3t_1t_2t_3$) is an odd integer, is
\begin{qoptions}
\opt{A}{$\dfrac12$}
\opt{B}{$\dfrac18$}
\opt{C}{$\dfrac1{16}$}
\opt{D}{$\dfrac1{32}$}
\end{qoptions}
\qn{5} If $A$ is an orthogonal matrix of order $3\times3$ and $\det(A)=1$ then which of the following matrices is/are \textbf{NOT} invertible
\begin{qoptions}
\opt{A}{$A-I$}
\opt{B}{$A^2-I$}
\opt{C}{$(A-A^T)^2$}
\opt{D}{$A-(\det(A))I$}
\end{qoptions}
\qn{6} Let for $n\in N$ ; $(2x^2-2x+1)^n=\displaystyle\sum_{r=0}^{2n}(-1)^ra_rx^r$ then
\begin{qoptions}
\opt{A}{$\displaystyle\sum_{r=0}^{2n}a_r=1$}
\opt{B}{$\displaystyle\sum_{r=1}^{2n}(-1)^ra_r=0$}
\opt{C}{$\displaystyle\sum_{k\ge0}{}^nC_k\,{}^{2n-2k}C_r=2^{n-r}a_r$}
\opt{D}{$\displaystyle\sum_{r=0}^{2n}\left((-1)^ra_ra_{2n-r}\right)=0$ if $n$ = odd}
\end{qoptions}
\qn{7} Let $(x-k)^2+y^2=r^2$, ($k,r>0$) is the largest circle, no part of which is lying outside the ellipse $3x^2+4y^2=12$. Now
\begin{qoptions}
\opt{A}{if $k<\dfrac12$ then $r+k=2$}
\opt{B}{if $k>\dfrac12$ then $r+k=2$}
\opt{C}{if $k<\dfrac12$ then $r+k<2$}
\opt{D}{if $k>\dfrac12$ then $r+k<2$}
\end{qoptions}
\sect{SEC--2}
\paraq{1}{3}
Consider the following path traced by point $(x,y)$ on the curve $x^a+y^a=1$ lying on the $xy$ plane.
\end{multicols}
\begin{center}\small\setlength{\tabcolsep}{4pt}
\begin{tabular}{|p{0.9cm}|p{3.6cm}|p{0.9cm}|p{3.6cm}|p{0.9cm}|p{3.6cm}|}\hline
\multicolumn{2}{|p{\dimexpr 4.5cm+2\tabcolsep\relax}|}{\centering \textbf{Column-I}\newline\textbf{(Condition)}} & \multicolumn{2}{p{\dimexpr 4.5cm+2\tabcolsep\relax}|}{\centering \textbf{Column-II}\newline\textbf{(Nature of path)}} & \multicolumn{2}{p{\dimexpr 4.5cm+2\tabcolsep\relax}|}{\centering \textbf{Column-III}\newline\textbf{(A tangent to the path)}}\\\hline
(I) & $a=-1$ & (i) & Circular & (P) & $x+y=\dfrac12$\\\hline
(II) & $a=\dfrac12$ & (ii) & Hyperbolic & (Q) & $x+y=1$\\\hline
(III) & $a=1$ & (iii) & Linear & (R) & $x+y=\sqrt2$\\\hline
(IV) & $a=2$ & (iv) & Parabolic & (S) & $x+y=4$\\\hline
\end{tabular}
\end{center}
\begin{multicols}{2}
\qn{1} Which of the following options is the only \textbf{CORRECT} combination?
\begin{qoptions}
\opt{A}{(III) (iii) (P)}
\opt{B}{(IV) (i) (S)}
\opt{C}{(II) (iv) (P)}
\opt{D}{(I) (ii) (R)}
\end{qoptions}
\qn{2} Which of the following options is the only \textbf{INCORRECT} combination ?
\begin{qoptions}
\opt{A}{(III) (iii) (Q)}
\opt{B}{(I) (ii) (S)}
\opt{C}{(III) (iii) (P)}
\opt{D}{(II) (iv) (P)}
\end{qoptions}
\qn{3} Which of the following options is the only \textbf{CORRECT} combination?
\begin{qoptions}
\opt{A}{(II) (iii) (P)}
\opt{B}{(I) (ii) (R)}
\opt{C}{(IV) (i) (Q)}
\opt{D}{(IV) (i) (R)}
\end{qoptions}
\paraq{4}{6}
Match the various functions $f(x)$ given in \textbf{Column-I} with their \textbf{RANGE} given in \textbf{Column-II} and a property of the real number $\alpha$ such that $f(\alpha)=1$, given in \textbf{Column-III}.\\
($\{\cdot\}$ and $[\cdot]$ respectively denote the fraction part and the greatest integer function)
\end{multicols}
\begin{center}\small\setlength{\tabcolsep}{4pt}
\begin{tabular}{|p{0.9cm}|p{3.6cm}|p{0.9cm}|p{3.6cm}|p{0.9cm}|p{3.6cm}|}\hline
\multicolumn{2}{|p{\dimexpr 4.5cm+2\tabcolsep\relax}|}{\centering \textbf{Column-I}\newline\textbf{(Function)}} & \multicolumn{2}{p{\dimexpr 4.5cm+2\tabcolsep\relax}|}{\centering \textbf{Column-II}\newline\textbf{(Range)}} & \multicolumn{2}{p{\dimexpr 4.5cm+2\tabcolsep\relax}|}{\centering \textbf{Column-III}\newline\textbf{($f(\alpha)=1$)}}\\\hline
(I) & $f(x)=[x]\,\{x\}$ & (i) & $(-\infty,\infty)$ & (P) & $\alpha_{\max.}=1$\\\hline
(II) & $f(x)=\bigl|[x]\bigr|\,\{x\}$ & (ii) & $[0,2)$ & (Q) & $|\alpha|_{\min.}=\dfrac32$\\\hline
(III) & $f(x)=\left[x\left[\dfrac1x\right]\right]$ & (iii) & $[0,\infty)$ & (R) & $|\alpha|_{\min.}=\dfrac52$\\\hline
(IV) & $f(x)=\{x\}+\left\{\dfrac1x\right\}$ & (iv) & $\{k:k\in I,\,k\ge0\}$ & (S) & $|\alpha-1|_{\min.}=\dfrac{\sqrt5-1}2$\\\hline
\end{tabular}
\end{center}
\begin{multicols}{2}
\qn{4} Which of the following options is the only \textbf{CORRECT} combination ?
\begin{qoptions}
\opt{A}{(I) (iii) (Q)}
\opt{B}{(I) (i) (R)}
\opt{C}{(III) (iv) (R)}
\opt{D}{(II) (iii) (S)}
\end{qoptions}
\qn{5} Which of the following options is the only \textbf{CORRECT} combination ?
\begin{qoptions}
\opt{A}{(I) (i) (P)}
\opt{B}{(IV) (iii) (S)}
\opt{C}{(III) (ii) (Q)}
\opt{D}{(II) (iii) (Q)}
\end{qoptions}
\qn{6} Which of the following options is the only \textbf{INCORRECT} combination ?
\begin{qoptions}
\opt{A}{(IV) (ii) (S)}
\opt{B}{(III) (iv) (P)}
\opt{C}{(III) (iv) (S)}
\opt{D}{(I) (i) (R)}
\end{qoptions}
\sect{SEC--3}
\qn{1} Let $\tan^{-1}x=\dfrac12\tan^{-1}y=\dfrac13\tan^{-1}z=\dfrac\pi7$, then the value of $\dfrac12\left(x^2+y^2+z^2-x^2y^2z^2\right)$, is
\qn{2} Let for the function $f(x)=5a^2x^3+6ax^2-27x+3b$, ($a\in R^-$, $b\in R$), $f'\left(\dfrac{-5}9\right)=0$, $f''\left(\dfrac{-5}9\right)<0$ and there exists some $\alpha\ne\dfrac{-5}9$ such that $f'(\alpha)=0$, $f(\alpha)>0$. The minimum possible integral value of $(a+b)$, is
\qn{3} Let \[I=\int_{-3}^{-2}\left(\dfrac{x^2-x}{x^3-3x+1}\right)^2\!\left(1+\dfrac1{x^2}+\dfrac1{(1-x)^2}\right)dx\] then the greatest integer less than the value of $I$, is
\qn{4} Let $f$ be a function such that $xf(x)=\pi\sin x\ \forall\ x\in(0,\pi)$ and $f(x)+f(\pi-x)=g(x)$ then $g_{\max.}$ is
\qn{5} Let $f(x)=(x-1)\,|\mathrm{sgn}(x)|$ and $g(x)=\max.(f(t))$, $-\infty<t\le x$, $-\infty<x<\infty$, then the number of points where $g(x)$ is non-differentiable is (sgn $(x)$ represents signum function of $x$)
\papersec{2}
\partsec{1}{PHYSICS}
\sect{SEC--1}
\qn{1} A thin non conducting disc of radius $R$ is rotating clockwise (see figure) with an angular velocity $\omega$ about its central axis, which is perpendicular to its plane. Both its surfaces carry +ve charges of uniform surface density. Half the disc is in a region of a uniform, unidirectional magnetic field $B$ parallel to the plane of the disc, as shown. Then,
\qfig{p2_phy_s1_q01_1}
\begin{qoptions}
\opt{A}{The net torque on the disc is zero.}
\opt{B}{The net torque vector on the disc is directed leftwards.}
\opt{C}{The net torque vector on the disc is directed rightwards.}
\opt{D}{The net torque vector on the disc is parallel to $B$.}
\end{qoptions}
\qn{2} At the initial moment three point $A$, $B$ and $C$ are on a horizontal plane along a straight line such that $AB=BC$. Point $A$ begins to move upward with a constant velocity `$v$' and point $C$ downward without any initial velocity at a constant acceleration `$a$'. If the points begin to move simultaneously, then the initial velocity and acceleration of point $B$ for all the three particles to be constantly on same straight line must be:
\begin{qoptions}
\opt{A}{$\dfrac v2$ upwards, $\dfrac a2$ downwards}
\opt{B}{$\dfrac v2$ upwards, $\dfrac a2$ upwards}
\opt{C}{$\dfrac v2$ downwards, $\dfrac a2$ downwards}
\opt{D}{$\dfrac v2$ downwards, $\dfrac a2$ upwards}
\end{qoptions}
\qn{3} The diagram shows three concentric conducting spherical shells having radii $R$, $2R$ and $3R$. The initial potential of each shell is as mentioned in the figure. If the inner most shell is earth then the charge present on its outer surface would be equal to
\qfig{p2_phy_s1_q03_1}
\begin{qoptions}
\opt{A}{zero}
\opt{B}{$40\pi\varepsilon_0R$}
\opt{C}{$4\pi\varepsilon_0R$}
\opt{D}{$-\pi\varepsilon_0R$}
\end{qoptions}
\qn{4} $ABCDECA$ is a planar body of mass $m$ of uniform thickness and same material. The dimensions are as shown in the figure. The moment of inertia of the body about an axis passing through point $A$ and perpendicular to planar body is $I_1$ and that of about an axis passing through $C$ and perpendicular to planer body is $I_2$. Then choose the correct relation :
\qfig{p2_phy_s1_q04_1}
\begin{qoptions}
\opt{A}{$\dfrac{I_1}{I_2}=3$}
\opt{B}{$\dfrac{I_1}{I_2}=\dfrac13$}
\opt{C}{$\dfrac{I_1}{I_2}=1$}
\opt{D}{$\dfrac{I_2}{I_1}=\dfrac23$}
\end{qoptions}
\qn{5} If an ellipse is rotated about $X$-axis an ellipsoid is formed, then the coordinates of the centre of mass of a solid uniform semi--ellipsoid formed by rotating a quarter ellipse given by the equation $\dfrac{x^2}{a^2}+\dfrac{y^2}{b^2}=1$ as shown in diagram is :-
\qfig{p2_phy_s1_q05_1}
\begin{qoptions}
\opt{A}{$\left(\dfrac{3a}8,0\right)$}
\opt{B}{$\left(\dfrac a8,0\right)$}
\opt{C}{$\left(\dfrac a4,0\right)$}
\opt{D}{$\left(\dfrac{2a}3,0\right)$}
\end{qoptions}
\qn{6} In a smooth stationary cart of length $d$, a small block is projected from the rear end of the cart along it's length with velocity $v$ towards front end. Coefficient of restitution for each collision is $e$. The cart rests on a smooth ground and can move freely. The time taken by block to come to rest w.r.t. cart is (Assume the time for collision is negligibly small).
\qfig{p2_phy_s1_q06_1}
\begin{qoptions}
\opt{A}{$\dfrac{ed}{(1-e)v}$}
\opt{B}{$\dfrac{ed}{(1+e)v}$}
\opt{C}{$\dfrac de$}
\opt{D}{Infinite}
\end{qoptions}
\qn{7} According to Moseley's law the ratio of the slopes of graph between $\sqrt f$ and $Z$ for $K_\beta$ and $K_\alpha$ is: (where the symbols have their usual meaning).
\begin{qoptions}
\opt{A}{$\sqrt{\dfrac{32}{27}}$}
\opt{B}{$\sqrt{\dfrac{27}{32}}$}
\opt{C}{$\sqrt{\dfrac{33}{22}}$}
\opt{D}{$\sqrt{\dfrac{22}{33}}$}
\end{qoptions}
\sect{SEC--2}
\qn{1} A micrometer has a resistance of 100$\Omega$ and a full scale range of 50$\mu$A. It can be used as a voltmeter or a higher range ammeter provided a resistance is added to it. Pick the correct range and resistance combination(s).
\begin{qoptions}
\opt{A}{50 V range with 10k$\Omega$ resistance in series}
\opt{B}{10 V range with 200k$\Omega$ resistance in series}
\opt{C}{5mA range with 1$\Omega$ resistance in parallel}
\opt{D}{10mA range with 1k$\Omega$ resistance in parallel}
\end{qoptions}
\qn{2} The two plates of an ideal capacitor of capacitance $C$ are given charges $Q_1$ and $Q_2$ as shown in the circuit below. This capacitor is connected across a resistance $R$ as shown key is closed at $t=0$. Then after time $t$.
\qfig{p2_phy_s2_q02_1}
\begin{qoptions}
\opt{A}{Total charge on the right plate is $q_2=\dfrac{Q_1+Q_2}2-\left(\dfrac{Q_1-Q_2}2\right)e^{-\frac t{RC}}$}
\opt{B}{Total charge on the left plate is $q=\dfrac{Q_1+Q_2}2+\left(\dfrac{Q_1-Q_2}2\right)e^{-\frac t{RC}}$}
\opt{C}{Initial potential difference across the plates is given by $\left|\dfrac{Q_1-Q_2}{2C}\right|$}
\opt{D}{Initial potential difference across the plates is given by $\dfrac{Q_1+Q_2}{2C}$}
\end{qoptions}
\qn{3} A circular conducting loop of radius $R$ and resistance per unit length $\lambda$ is pulled out from the region of uniform magnetic field with constant velocity $v$. The situation shown in the figure corresponds to time $t=0$. Mark out the correct statement(s)
\qfig{p2_phy_s2_q03_1}
\begin{qoptions}
\opt{A}{Just after $t=0$ i.e., the motion starts, the induced current in the loop is $\dfrac{\sqrt3Bv}{2\pi\lambda}$}
\opt{B}{Current will be induced in the loop for $\dfrac{3R}{2v}$ sec}
\opt{C}{At any time $t$, the current induced in the loop is given by $\dfrac{Bv\sqrt{3R^2-4v^2t^2+4Rvt}}{\lambda\times2\pi R}$ and is in clockwise direction.}
\opt{D}{Induced current is in clockwise direction for $t=0$ to $\dfrac R{2v}$ and thereafter it becomes in anticlockwise direction.}
\end{qoptions}
\qn{4} Two AC sources $\xi_1=10\sin(\omega t)$V and $\xi_2=10\sin\left(\omega t-\dfrac\pi3\right)$V are fed to an $LR$ circuit. The resistance is $R=10\ \Omega$ and inductive reactance is $X_L=10\sqrt3\Omega$ :-
\qfig{p2_phy_s2_q04_1}
\begin{qoptions}
\opt{A}{The equivalent emf of the AC source has a peak value of $10\sqrt3$V.}
\opt{B}{The peak value of current is $\dfrac{\sqrt3}2$A}
\opt{C}{The voltage across the inductor at $t=\dfrac\pi{3\omega}$ is $\dfrac{15\sqrt3}2$V}
\opt{D}{The voltage across the resistor at $t=\dfrac\pi{3\omega}$ is $\dfrac{15}2$V}
\end{qoptions}
\qn{5} A positively charged particle of charge $q$ and mass $m$ is suspended from a point by a string of length $\ell$ in a vertical plane. In the entire space a uniform horizontal electric field $E$ exists. The particle is drawn aside so that string becomes vertical and then it is projected horizontally with velocity $v$ such that particle starts to move along a circle with some constant speed $v$. Then choose the correct statements.
\qfig{p2_phy_s2_q05_1}
\begin{qoptions}
\opt{A}{Speed of the particle is $v=\dfrac{qE}m\sqrt{\dfrac\ell g}$}
\opt{B}{Tension in the string is $mg+\dfrac{q^2E^2}{mg}$}
\opt{C}{Time period of revolution is $2\pi\sqrt{\dfrac{\ell g}{g^2+\left(\frac{qE}m\right)^2}}$}
\opt{D}{Time period of revolution is independent of mass of particle.}
\end{qoptions}
\qn{6} In the figure shown, the mass of the disc as well as that of the trolley is $M$. The spring is ideal and has stiffness $k$. The trolley can move horizontally on the floor without friction, and the disc can roll on the trolley surface without slipping. The spring is compressed and the system released so that oscillations begin. At the instant of release, magnitude of :
\qfig{p2_phy_s2_q06_1}
\begin{qoptions}
\opt{A}{Acceleration of center of disc = twice of that of trolley}
\opt{B}{Acceleration of center of disc = thrice of that of trolley}
\opt{C}{Spring force is four times of friction.}
\opt{D}{Spring force is two times of friction.}
\end{qoptions}
\qn{7} An isotropic sound source $A$ is moving in a circle of radius $R$ with small speed $v$. An observer $B$ is hearing this sound (see figure).
\qfig{p2_phy_s2_q07_1}
\begin{qoptions}
\opt{A}{The intensity of sound heard by $B$ will be maximum just after $A$ crosses 1.}
\opt{B}{The frequency of sound heard by $B$ will be minimum when $A$ is at 6.}
\opt{C}{The frequency of sound heard by $B$ will be minimum just after $A$ crosses 2.}
\opt{D}{The frequency heard by $B$ when source is at 1 \& 4 is minimum.}
\end{qoptions}
\sect{SEC--3}
\paraq{1}{2}
In YDSE arrangement shown in figure, fringes are seen on screen using monochromatic source $S$ having wavelength 3000\AA\ (in air). $S_1$ and $S_2$ are two slits separated by $d=1$mm from each other and $D=1$m. Left of slits $S_1$ and $S_2$ medium of refractive index $\mu_1=2$ is present \& to the right of slits $S_1$ and $S_2$ medium of refractive index $\mu_2=\dfrac32$ is present. A thin slab of thickness `$t$' is placed in front of $S_1$. The refractive index $\mu_3$ of slab varies with distance ($x$) from its starting face as shown in figure.
\qfig{p2_phy_s3_para1_2_1}
\qn{1} In order to get central maxima at the centre of the screen ($O$), the thickness of slab required is :-
\begin{qoptions}
\opt{A}{1$\mu$m}
\opt{B}{2$\mu$m}
\opt{C}{0.5$\mu$m}
\opt{D}{1.5 $\mu$m}
\end{qoptions}
\qn{2} If the thickness of the slab is selected 1mm, then the position of the central maxima will be ($Y$ - coordinate) Taking $O$ to be origin
\begin{qoptions}
\opt{A}{$\dfrac13$mm}
\opt{B}{$-\dfrac13$mm}
\opt{C}{$\dfrac16$mm}
\opt{D}{$-\dfrac16$mm}
\end{qoptions}
\paraq{3}{4}
Air Standard Otto cycle :\\
The air-standard-Otto cycle is the ideallzed cycle for the spark-ignition internal combustion engines. This cycle is shown on P-V and T-S diagrams. The Otto cycle 1-2-3-4 consists of following four process:\\
Process 1-2: Reversible adiabatic compression of air.\\
Process 2-3: Heat addition at constant volume.\\
Process 3-4: Reversible adiabatic expansion of air.\\
Process 4-1: Heat rejection at constant volume.
\qfig{p2_phy_s3_para3_4_1}
Mean Effective Pressure is defined as the ratio of the net work done to the displacement volume of the piston.
\qn{3} For a given mass of working gas of adiabatic constant $\gamma\left(=\dfrac{C_P}{C_V}\right)$ if the compression ratio $r$ is defined as $r=\dfrac{V_1}{V_2}$ (Where $V_1$ is volume of gas at 1 and $V_2$ is volume of gas at 2). Then the efficiency of the cycle is given as.
\begin{qoptions}
\opt{A}{$n_{th}=1-\left(\dfrac1r\right)^{\gamma-1}$}
\opt{B}{$n_{th}=1+\left(\dfrac1r\right)^{\gamma-1}$}
\opt{C}{$n_{th}=1+\left(\dfrac1r\right)^{\gamma+1}$}
\opt{D}{$n_{th}=1-\left(\dfrac1r\right)^{\gamma+1}$}
\end{qoptions}
\qn{4} If the displacement volume is given by ($V_1-V_2$), then the mean effective pressure of the above cycle is given as
\begin{qoptions}
\opt{A}{$\left(\dfrac1{\gamma-1}\right)\left(\dfrac{P_1}{T_1}\right)\left(\dfrac r{r-1}\right)\{(T_3-T_2)-(T_4-T_1)\}$}
\opt{B}{$\left(\dfrac1{\gamma+1}\right)\left(\dfrac{P_1}{T_1}\right)\left(\dfrac r{r-1}\right)\{(T_3-T_2)+(T_4-T_1)\}$}
\opt{C}{$\left(\dfrac1{\gamma-1}\right)\left(\dfrac{P_1}{T_1}\right)\left(\dfrac r{r+1}\right)\{(T_3-T_2)+(T_4-T_1)\}$}
\opt{D}{$\left(\dfrac1{\gamma-1}\right)\left(\dfrac{P_1}{T_1}\right)\left(\dfrac r{r-1}\right)\{(T_3+T_2)+(T_4-T_1)\}$}
\end{qoptions}
\partsec{2}{CHEMISTRY}
\sect{SEC--1}
\qn{1} Choose correct Statement:
\begin{qoptions}
\opt{A}{In group 15 Hydrides, $\mathrm{NH_3}$ is least volatile.}
\opt{B}{$\mathrm{HNO_2}$ can act as oxidiser as well as reducer.}
\opt{C}{$\mathrm{NH_3}$ is a non-polar molecule.}
\opt{D}{$\mathrm{N_2}$ is a reactive gas at room temperature.}
\end{qoptions}
\qn{2} Choose incorrect statement :
\begin{qoptions}
\opt{A}{$\mathrm{SnF_4}$ is more ionic than $\mathrm{SnI_4}$.}
\opt{B}{Cement contains silicates.}
\opt{C}{Diamond does not conduct electricity.}
\opt{D}{$\mathrm{CCl_4}$ easily hydrolyses to give $\mathrm{CH_4}$ at room temperature.}
\end{qoptions}
\qn{3} In given reactions which one is correct ?
\begin{qoptions}
\opt{A}{\qfigi{p2_chem_s1_q03_1}}
\opt{B}{\qfigi{p2_chem_s1_q03_B}}
\opt{C}{\qfigi{p2_chem_s1_q03_C}}
\opt{D}{\qfigi{p2_chem_s1_q03_D}}
\end{qoptions}
\qn{4} Which compound will not give positive $\mathrm{NaHCO_3}$ test ?
\begin{qoptions}
\opt{A}{Picric acid}
\opt{B}{Benzoic acid}
\opt{C}{Carbinol}
\opt{D}{Terephthalic acid}
\end{qoptions}
\qn{5} Identify reaction that correctly match with their major product.
\begin{qoptions}
\opt{A}{\qfigi{p2_chem_s1_q05_A}}
\opt{B}{\qfigi{p2_chem_s1_q05_B}}
\opt{C}{\qfigi{p2_chem_s1_q05_C}}
\opt{D}{\qfigi{p2_chem_s1_q05_D}}
\end{qoptions}
\qn{6} Which of the following statements is correct ?
\begin{qoptions}
\opt{A}{Milk is not an emulsion.}
\opt{B}{Lyophilic hydrosols have lower surface tension than $\mathrm{H_2O}$.}
\opt{C}{Lyophobhic hydrosols have lower surface tension than water.}
\opt{D}{Conductivity of Lyophilic colloids is negligible}
\end{qoptions}
\qn{7} Bombardment on nucleus $A$ by $\alpha$-particle leads to its artificial disintegration in two ways, (i) and (ii) as shown. Products $X$, $Y$ and $Z$ respectively are :
\qfig{p2_chem_s1_q07_1}
\begin{qoptions}
\opt{A}{proton, neutron, positron}
\opt{B}{neutron, positron, proton}
\opt{C}{proton, positron, neutron}
\opt{D}{positron, proton, neutron}
\end{qoptions}
\sect{SEC--2}
\qn{1} Which carbonyl compound(s) give racemic cyanohydrin, when react with HCN.
\begin{qoptions}
\opt{A}{Acetone}
\opt{B}{Cyclopentanone}
\opt{C}{Equal amount of ($\pm$) \qfigi{p2_chem_s2_q01_1}}
\opt{D}{Acetaldehyde}
\end{qoptions}
\qn{2} In given reaction, which one is/are correct according to major product.
\begin{qoptions}
\opt{A}{\qfigi{p2_chem_s2_q02_A}}
\opt{B}{\qfigi{p2_chem_s2_q02_B}}
\opt{C}{\qfigi{p2_chem_s2_q02_C}}
\opt{D}{\qfigi{p2_chem_s2_q02_D}}
\end{qoptions}
\qn{3} Choose correct statement(s)
\begin{qoptions}
\opt{A}{Hydrolysis of $\mathrm{Mg_2C_3}$ involves redox reaction.}
\opt{B}{$\mathrm{NF_3}$ does not undergo hydrolysis in ordinary conditions.}
\opt{C}{$\mathrm{B_3N_3H_6}$ produces no gaseous product when hydrolysed at high temperature}
\opt{D}{The products of hydrolysis of hypophosphoric acid ($\mathrm{H_4P_2O_6}$) are $\mathrm{H_3PO_3}$ and $\mathrm{H_3PO_4}$.}
\end{qoptions}
\qn{4} Which of the following reaction(s) will produce at least one gaseous compound ?
\begin{qoptions}
\opt{A}{$\mathrm{Na_2CO_3(aq)+dil.H_2SO_4}\longrightarrow$}
\opt{B}{$\mathrm{Na_2SO_3(aq)+dil.H_2SO_4}\longrightarrow$}
\opt{C}{$\mathrm{Na_2C_2O_4(aq)+conc.H_2SO_4}\longrightarrow$}
\opt{D}{$\mathrm{NaNO_3(aq)+conc.H_2SO_4}\xrightarrow{\Delta}$}
\end{qoptions}
\qn{5} Anhydrous $\mathrm{CaCl_2}$ is often used as desiccant. In the presence of excess $\mathrm{CaCl_2}$ the amount of water taken up is governed by $K_P=6.4\times10^{85}$ for the following reaction at room temperature.\\
$\mathrm{CaCl_2(s)+6H_2O(g)\rightleftharpoons CaCl_2{\cdot}6H_2O(s)}$\\
Which of the following statements is correct about above reaction ?
\begin{qoptions}
\opt{A}{At equilibrium addition of $\mathrm{CaCl_2}$ will shift reaction in forward direction.}
\opt{B}{Vapour pressure of water at equilibrium in a closed vessel that contains $\mathrm{CaCl_2}$(s) is $5\times10^{-15}$ atm approximately.}
\opt{C}{Vapour pressure of water at equilibrium in a closed vessel becomes half if volume of container is reduced to half at same temperature.}
\opt{D}{If temperature is increased, then vapour pressure of water at equilibrium in closed vessel decreases.}
\end{qoptions}
\qn{6} Which of the following statements is/are correct ?
\begin{qoptions}
\opt{A}{At constant volume, there will be no effect of changing $T$ or $P$ on mean free path.}
\opt{B}{Collision number increases on increasing temperature at constant volume.}
\opt{C}{If diffusion coefficient is proportional to mean free path and mean speed then on increasing absolute temperature of an ideal gas by four times and pressure by two times, diffusion coefficient of gas increases four times.}
\opt{D}{According to kinetic theory of gases, between two successive colloisions, a gas molecule travels in straight line path.}
\end{qoptions}
\qn{7} Which reaction is/are correct ?
\begin{qoptions}
\opt{A}{\qfigi{p2_chem_s2_q07_A}}
\opt{B}{\qfigi{p2_chem_s2_q07_B}}
\opt{C}{\qfigi{p2_chem_s2_q07_C}}
\opt{D}{\qfigi{p2_chem_s2_q07_D}}
\end{qoptions}
\sect{SEC--3}
\paraq{1}{2}
d-Block metals are extracted from their ores by a wide variety of techniques. Most metals are obtained by direct treatment of their ores with chemical reagents, but the extraction of certain metals require electrolysis . Pyrolusite ($\mathrm{MnO_2}$) Hausmannite ($\mathrm{Mn_3O_4}$) Rhodocroside ($\mathrm{MnCO_3}$) are the natural abundent ore of Mn use to extract metal and pyrolusite ($\mathrm{MnO_2}$) ore is used to prepare some compounds, such as $\mathrm{KMnO_4}$ can be prepared with intermediate fomation of $\mathrm{K_2MnO_4}$.
\qn{1} The set of metals whose extraction involves hydrometallurgy is:
\begin{qoptions}
\opt{A}{Al, Zn}
\opt{B}{Fe, Cu}
\opt{C}{Ag, Au}
\opt{D}{Fe, Ag}
\end{qoptions}
\qn{2} Consider following Statements\\
I: $\mathrm{KMnO_4}$ can act as oxidiser in acidic, basic and neutral medium.\\
II: $\mathrm{MnO_4^{2-}}$ is paramagnetic while $\mathrm{MnO_4^{-1}}$ is diamagnetic.\\
III. $\mathrm{MnO_4^-}$ is only oxidising agent but $\mathrm{MnO_4^{2-}}$ can act as oxidising as well as reducing agent\\
IV. $\mathrm{Mn_3O_4}$ can be converted to Mn by aluminothermite process\\
The Correct statement is/are:
\begin{qoptions}
\opt{A}{I, II \& III only}
\opt{B}{I, II \& IV only}
\opt{C}{I, II, III \& IV}
\opt{D}{I \& IV only}
\end{qoptions}
\paraq{3}{4}
The electronic ground state of hydrogen atom contains one electron in first orbit. If sufficient energy is provided, the electron can be promoted to higher energy level. The electronic energy of hydrogen like species (any atom/ions with nuclear charge $|Ze|$ and one electron) can be given as
\[E_n=-R_H\,\frac{Z^2}{n^2}\]
Where\\
$R_H$ = Rydberg's constant\\
$n$ = Principle quantum number , $Z$ = Atomic number
\qn{3} The ratio of energy of an electron in the ground state of $\mathrm{Be^{3+}}$ ion to that of ground state of H-atom is
\begin{qoptions}
\opt{A}{16}
\opt{B}{4}
\opt{C}{1}
\opt{D}{8}
\end{qoptions}
\qn{4} The experimental I.E.$_1$ of He atom is 24.58 eV. Total energy required to remove two electrons from He is
\begin{qoptions}
\opt{A}{49.16 eV}
\opt{B}{78.98 eV}
\opt{C}{29.82 eV}
\opt{D}{None of these}
\end{qoptions}
\partsec{3}{MATHEMATICS}
\sect{SEC--1}
\qn{1} If line $x=y=z$ intersect the line $x\sin A+y\sin B+z\sin C-2d^2=0=x\sin2A+y\sin2B+z\sin2C-d^2$ then $\sin\dfrac A2\sin\dfrac B2\sin\dfrac C2$ equals, where $(A+B+C)=\pi$
\begin{qoptions}
\opt{A}{$\dfrac1{16}$}
\opt{B}{$\dfrac18$}
\opt{C}{$\dfrac1{32}$}
\opt{D}{$\dfrac1{12}$}
\end{qoptions}
\qn{2} Solution of $\dfrac{dy}{dx}=\dfrac{y^3}{e^{2x}+y^2}$ is
\begin{qoptions}
\opt{A}{$e^{-2x}y^2+2\ln|y|=C$}
\opt{B}{$e^{2x}y^2-2\ln|y|=C$}
\opt{C}{$e^x+\ln|y|=C$}
\opt{D}{None of these}
\end{qoptions}
\qn{3} Let $f(x)=x^2-2x$, $g(x)=f(f(x)-1)+f(5-f(x))$ then
\begin{qoptions}
\opt{A}{$g(x)<0\ \forall\ x\in R$}
\opt{B}{$g(x)<0$ for some $x\in R$}
\opt{C}{nothing can be said about $g(x)$}
\opt{D}{$g(x)\ge0\ \forall\ x\in R$}
\end{qoptions}
\qn{4} If $[\bar a\ \bar b\ \bar c]=|\bar d|=6$, then $\bigl|[\bar d\ \bar b\ \bar c]\,\bar a+[\bar d\ \bar c\ \bar a]\,\bar b+[\bar d\ \bar a\ \bar b]\,\bar c\bigr|=$
\begin{qoptions}
\opt{A}{36}
\opt{B}{12}
\opt{C}{24}
\opt{D}{6}
\end{qoptions}
\qn{5} Line segments drawn from vertex opposite to the hypotenuse of a right angle triangle to the points trisecting the hypotenuse, have lengths $\sin x$ and $\cos x$ $\left(0<x<\dfrac\pi2\right)$. Then length of hypotenuse is
\begin{qoptions}
\opt{A}{$\dfrac13$}
\opt{B}{$\dfrac32$}
\opt{C}{$\dfrac{3\sqrt5}5$}
\opt{D}{can't be determined}
\end{qoptions}
\qn{6} Let $x+y=0$ and $x-y=0$ are tangents drawn to a parabola whose focus is at $S(1,2)$ then length of latus rectum of parabola is
\begin{qoptions}
\opt{A}{$\dfrac6{\sqrt5}$}
\opt{B}{$\dfrac4{\sqrt5}$}
\opt{C}{$\dfrac2{\sqrt5}$}
\opt{D}{$\dfrac1{\sqrt5}$}
\end{qoptions}
\qn{7} If $A$ is an involutary matrix and diagonal matrix of non positive entries and order 3 then
\begin{qoptions}
\opt{A}{there may exist some diagonal element in $A$, which is/are zero}
\opt{B}{$|A|=-1$}
\opt{C}{$A^{-1}$ does not exist}
\opt{D}{$|3\,\mathrm{adj}.2A|=-1728$}
\end{qoptions}
\sect{SEC--2}
\qn{1} If $4a^2-5b^2+6a+1=0$ and the line $ax+by+1=0$ touches a fixed circle, then
\begin{qoptions}
\opt{A}{centre of circle is $(3,0)$}
\opt{B}{radius of circle $=\sqrt5$}
\opt{C}{circle passes through $(0,1)$}
\opt{D}{None of these}
\end{qoptions}
\qn{2} Number of real solution(s) of $2^{\sin^{-1}x}+2^{\cos^{-1}x}+2^{\tan^{-1}x}=3$ is less than
\begin{qoptions}
\opt{A}{0}
\opt{B}{1}
\opt{C}{2}
\opt{D}{3}
\end{qoptions}
\qn{3} Let $\displaystyle\int_0^x\dfrac{(m t\cos4t-\ell\sin4t)\,dt}{t^2}=\dfrac{\ell\sin4x}x-1$, where $\left(0<x<\dfrac\pi4\right)$, then
\begin{qoptions}
\opt{A}{$\ell+m=\dfrac54$}
\opt{B}{$m=4\ell$}
\opt{C}{$\ell-m=\dfrac34$}
\opt{D}{$m=\ell+\dfrac14$}
\end{qoptions}
\qn{4} If tangent of any member of family of hyperbola $xy=4\sin^2\theta$, $\theta\in(0,2\pi)-\{\pi\}$ is not a normal to member of family of circles $x^2+y^2-2x-2y+\mu=0$, $\mu$ is real parameter, then $\theta$ belongs to
\begin{qoptions}
\opt{A}{$\left(\dfrac{5\pi}6,\dfrac{7\pi}6\right)$}
\opt{B}{$\left(0,\dfrac\pi6\right)$}
\opt{C}{$\left(\dfrac{11\pi}6,2\pi\right)$}
\opt{D}{$\left(\dfrac\pi6,\dfrac{5\pi}6\right)$}
\end{qoptions}
\qn{5} $f(x)=[\tan^{100}x]$ where $[\cdot]$ greatest integer function. Then
\begin{qoptions}
\opt{A}{$f(x)$ is discontinuous at $x=\tan^{-1}(3^{1/100})$, $\tan^{-1}(100^{1/100})$}
\opt{B}{$f(x)$ discontinuous at $x=\tan^{-1}(3^{1/100})$, $\tan^{-1}(99^{1/100})$}
\opt{C}{$\displaystyle\int_0^{\tan^{-1}(4^{1/100})}f(x)dx$\\$=3\tan^{-1}(4^{1/100})-\tan^{-1}(3^{1/100})$\\$\quad-\tan^{-1}(2^{1/100})-\dfrac\pi4$}
\opt{D}{$\displaystyle\int_0^{\tan^{-1}(4^{1/100})}f(x)dx$\\$=3\tan^{-1}4^{1/100}-2\tan^{-1}3^{1/100}$}
\end{qoptions}
\qn{6} Area of region given by curves $1-y^2\ge|x|$ and $|x|+|y|\ge1$ is $A$ and area of region $\{(x,y):x^2\le y\le|x|\}$ is $B$. Then $\dfrac AB$ is less than
\begin{qoptions}
\opt{A}{1}
\opt{B}{2}
\opt{C}{3}
\opt{D}{4}
\end{qoptions}
\qn{7} Set of all real values of $p$ such that both the roots of the equation $(p-5)x^2-2px+(p-4)=0$ are positive, one is less than 2 and other is lying between 2 and 3 is
\begin{qoptions}
\opt{A}{$\left(\dfrac{49}4,24\right)$}
\opt{B}{$(5,\infty)$}
\opt{C}{$(-\infty,4)\cup\left(\dfrac{49}4,\infty\right)$}
\opt{D}{$\left(4,\dfrac{49}4\right)$}
\end{qoptions}
\sect{SEC--3}
\paraq{1}{2}
Let $A=\{1,2,3,\ldots..n\}$ is a set containing $n$ elements.
\qn{1} For any given $n$, number of subsets of $A$ having $k(k<n)$ as largest element, is
\begin{qoptions}
\opt{A}{$2^k$}
\opt{B}{$n-k$}
\opt{C}{$2^{k-1}$}
\opt{D}{None of these}
\end{qoptions}
\qn{2} The number of subsets of $A$ having at least one odd integer is\\
{[}where [.] represents greatest integer function{]}
\begin{qoptions}
\opt{A}{$2^{\left[\frac n2\right]}$}
\opt{B}{$2^{\left[\frac{n+1}2\right]}$}
\opt{C}{$2^n-2^{\left[\frac n2\right]}$}
\opt{D}{$2^{\left[\frac n2\right]+1}$}
\end{qoptions}
\paraq{3}{4}
Let $S$ be the set of points $\{(x,y):0\le x\le3,\ 0\le y\le4$ where $x,y\in I\}$ two distinct points are randomly chosen from $S$.
\qn{3} The probability that mid point of the line segment of two points chosen belongs to $S$ is
\begin{qoptions}
\opt{A}{$\dfrac{21}{95}$}
\opt{B}{$\dfrac{42}{95}$}
\opt{C}{$\dfrac{36}{95}$}
\opt{D}{None of these}
\end{qoptions}
\qn{4} If two points $(x_1,y_1)$ and $(x_2,y_2)$ are chosen such that $x_1=x_2$ and $y_1\ne y_2$ then probability that mid point of these two points belongs to $S$, is
\begin{qoptions}
\opt{A}{$\dfrac15$}
\opt{B}{$\dfrac25$}
\opt{C}{$\dfrac8{15}$}
\opt{D}{$\dfrac8{95}$}
\end{qoptions}

\end{multicols}
\end{document}
